\subsection{SF-E808. Odd radicands ramified at two: an unrestricted
Chern-class
obstruction}\label{SF-E808}

Let \(F=\mathbf Q(\sqrt{\Delta})\), with \(\Delta>1\) squarefree and
\(\Delta\equiv3\pmod4\). Suppose a polarized abelian fourfold
\((X,\theta)\) has Rosati-invariant RM by \(F\), with tangent
multiplicities \((2,2)\), and that \(L=H_1(X,\mathbf Z_2)\) is stable
under the maximal local order \(R=\mathbf Z_2[u]\), \(u^2=\Delta\).
Suppose \(D=\int_X\theta^4/24\) is odd. Define

\[
 C=\frac{\theta_+^2+\theta_-^2}{2}.
\]

\textbf{Assertion.} For every integral class
\(\ell\in H^2(X,\mathbf Z_2)\), every positive integer \(r\), and every
two-adic unit \(z\in\mathbf Q_2\),

\[
 \frac{\ell^2}{2r}-zC\notin H^4(X,\mathbf Z_2).
 \tag{E808.1}
\]

Consequently no positive-rank perfect object, with any rational divisor
normalization, has normalized balanced character \(r v_q\) when \(rq\)
is a two-adic unit. A simple gluable fully semiregular lift with an
effective finite coherent target action \(G\) as in supplied S-22
therefore satisfies \(2\mid D|G|\).

Together with SF-E807, this removes the minimal-divisor-space hypothesis
from supplied S-32 for all real quadratic fields ramified at two.
Maximal-order stability at two and oddness of the relevant degree
product are retained.

\subsubsection{The local polarized RM
lattice}\label{the-local-polarized-rm-lattice}

The element \(u-1\) satisfies

\[
 t^2+2t+1-\Delta=0.
\]

This polynomial is Eisenstein at two, because \(1-\Delta\) has valuation
one. Thus \(R\) is a discrete valuation ring with uniformizer
\(\pi=u-1\). Concretely \(R=\mathbf Z_2+\mathbf Z_2\pi\), and the
valuation of \(a+b\pi\) is \(\min\{2v_2(a),2v_2(b)+1\}\); the two terms
have distinct valuation parity. This identifies \(R\) with the valuation
ring. The finite torsion-free \(R\)-module \(L\) is free of rank four.

Let \(E\) be the polarization pairing on \(L\). Its determinant is odd,
so it is perfect over \(\mathbf Z_2\). Exactly as in the full local
proof of SF-E807, the form

\[
 \Psi(x,y)=E(x,uy)+uE(x,y)
\]

is alternating and \(R\)-bilinear, and it is perfect because

\[
 E(x,y)=\operatorname{Tr}_{(F\otimes\mathbf Q_2)/\mathbf Q_2}
              \left(\frac{\Psi(x,y)}{2u}\right).
\]

For completeness, the trace pairing on \(R\) is perfect: it sends
\(a+bu\) to \(b\), and multiplication first by \(u\) sends it to \(a\).
This identifies the \(R\)-linear dual with the \(\mathbf Z_2\)-linear
dual. A symplectic \(R\)-basis \(x_i,y_i\), \(i=1,2\), follows by
splitting off pairs with pairing one.

Let \(\alpha_i,\beta_i,\gamma_i,\delta_i\) be the dual integral basis to
\(x_i,ux_i,y_i,uy_i\). Put

\[
 \begin{aligned}
 S&=\langle\alpha_1,\gamma_1,\alpha_2,\gamma_2\rangle,
 &T&=\langle\beta_1,\delta_1,\beta_2,\delta_2\rangle,\\
 \omega&=\alpha_1\gamma_1+\alpha_2\gamma_2,
 &\tau&=\beta_1\delta_1+\beta_2\delta_2,\\
 m&=\alpha_1\delta_1+\beta_1\gamma_1
       +\alpha_2\delta_2+\beta_2\gamma_2.
 \end{aligned}
\]

Evaluation of the pairings gives \(\theta=m\) and
\(\theta_u=\omega+\Delta\tau\). Hence

\[
 Q:=2C=m^{[2]}+\frac1\Delta\omega^{[2]}
            +\omega\tau+\Delta\tau^{[2]}.
 \tag{E808.2}
\]

This is an integral two-adic four-form. Its reduction is independent of
the odd radicand:

\[
 \bar Q=m^{[2]}+\omega^{[2]}+\omega\tau+\tau^{[2]}.
 \tag{E808.3}
\]

Here the displayed operations are divided powers and exterior products
reduced from the integral lattice; signs disappear modulo two.

\subsubsection{A contraction with no wedge
annihilator}\label{a-contraction-with-no-wedge-annihilator}

Write \(s_0=\alpha_1,s_1=\gamma_1,s_2=\alpha_2,s_3=\gamma_2\), and
\(t_0=\beta_1,t_1=\delta_1,t_2=\beta_2,t_3=\delta_2\). Let \(\lambda\)
be the dual functional selecting \(s_0\). Direct contraction of (E808.3)
gives

\[
 \chi=\iota_\lambda\bar Q
 =s_1s_2s_3+s_1t_2t_3+s_2t_1t_3+s_3t_1t_2.
 \tag{E808.4}
\]

The map

\[
 \bar L^*\longrightarrow\bigwedge^4\bar L^*,\qquad
 v\longmapsto v\chi
 \tag{E808.5}
\]

is injective. This is a direct monomial calculation, not a generic-rank
assertion. Rename the six variables appearing in \(\chi\) as
\(a,b,c,d,e,f\), so \(\chi=abc+aef+bdf+cde\). Their six exterior
products are

\[
 \begin{array}{c|l}
 a&abdf+acde\\
 b&abef+bcde\\
 c&acef+bcdf\\
 d&abcd+adef\\
 e&abce+bdef\\
 f&abcf+cdef .
 \end{array}
\]

The twelve monomials are distinct, proving independence of these six
images. Exterior multiplication by either of the remaining variables
\(s_0,t_0\) gives a nonzero form containing that variable. These two
images and the preceding six belong to separate summands according to
their \(s_0,t_0\) content. This proves injectivity on the entire
eight-dimensional space.

In particular \(\bar Q\ne0\), so \(Q\) is primitive at two.

\subsubsection{A divided-square lemma with denominator
two}\label{a-divided-square-lemma-with-denominator-two}

Let \(M\) be a free \(\mathbf Z_2\)-module, let \(Q\in\bigwedge^4M\),
and suppose some functional \(\lambda\in(M/2M)^*\) has the property that
exterior multiplication by \(\chi=\iota_\lambda\bar Q\) is injective on
\(M/2M\). Then there are no \(\ell\in\bigwedge^2M\) and positive integer
\(r\) with

\[
 \ell^{[2]}-\frac r2Q\in r\bigwedge^4M.
 \tag{E808.6}
\]

\textbf{Proof.} The hypothesis makes \(Q\) primitive, so integrality
forces \(r\) even. Write \(r=2^e k\), with \(k\) odd and \(e\ge1\). The
alleged congruence is

\[
 \ell^{[2]}\equiv 2^{e-1}kQ\pmod{2^e}.
 \tag{E808.7}
\]

The minimum valuation of its left side is exactly \(e-1\). Removing the
largest common factor \(2^a\) from \(\ell\) reduces the congruence to
the same one for a primitive two-form and exponent \(e'=e-2a\ge1\). We
may assume that \(\ell\) is primitive.

For \(e=1\), the congruence says \(\bar\ell^{[2]}=\bar Q\). The
contraction identity for divided squares gives

\[
 \chi=(\iota_\lambda\bar\ell)\bar\ell.
\]

If \(\iota_\lambda\bar\ell=0\), this contradicts \(\chi\ne0\). If it is
nonzero, it wedges to zero with \(\chi\), contradicting the assumed
injectivity. Thus \(e=1\) is impossible.

For \(e\ge2\), the alternating elementary-divisor algorithm gives

\[
 \ell=\alpha+2^{e-1}\beta,
\]

where \(\alpha\) is primitive and decomposable. Indeed primitivity makes
the first elementary divisor a unit, while the valuation of the divided
square makes the second elementary divisor have valuation \(e-1\); all
subsequent divisors have at least that valuation. Dividing (E808.7) by
\(2^{e-1}\) and reducing modulo two gives

\[
 \bar Q=\bar\alpha\bar\beta.
\]

The two-dimensional support \(P\) of \(\bar\alpha\) therefore wedges to
zero with \(\bar Q\). There is a nonzero \(v\in P\cap\ker\lambda\).
Contracting \(v\bar Q=0\) by \(\lambda\) gives

\[
 0=\lambda(v)\bar Q+v\iota_\lambda\bar Q=v\chi,
\]

which again contradicts injectivity. Both cases are impossible. ∎

\subsubsection{Chern-character consequence and
scope}\label{chern-character-consequence-and-scope}

Multiplication by the unit \(z\) in (E808.1) leaves \(\bar Q\) and its
contraction unchanged. The lemma applied to \(zQ\), with \(Q=2C\),
proves (E808.1). For a normalized balanced character, degree two forces
\(B=-c_1(E)/r\), while degree four gives

\[
 \operatorname{ch}_2(E)=\frac{c_1(E)^2}{2r}-(rq)C.
\]

The Chern character on the left is integral on a complex torus, since
\(\operatorname{ch}_2=c_1^2/2-c_2\) and every integral two-form has an
integral divided square. Hence \(rq\) cannot be a two-adic unit.

If a simple gluable fully semiregular equivariant lift had \(D|G|\) odd,
supplied S-22 would give \(4D(rq)^2/|G|\in\{2,4\}\). The value two has
impossible odd valuation for a rational square; the value four makes
\(rq\) a two-adic unit. Thus \(2\mid D|G|\).

No restriction to the RM divisor plane was made: every integral
\(H^2\)-class is included, so extra CM divisor classes cannot evade the
conclusion. This does not exclude even-degree polarizations, even-order
effective finite actions, nonmaximal local orders, or other proposed
characters. It constructs no eightfold supplier.

\subsubsection{Source and check status}\label{source-and-check-status}

The proof uses only the supplied S-22 Euler restriction for the final
semiregularity corollary. The integral obstruction, local pairing
calculation, contraction matrix, and lifting argument are proved in full
above. S-32 is a comparison result, not an input to this proof. The
finite certificate checks exact normal-form substitutions and the
contraction rank; it does not replace the proof for every rank or every
radicand.
